\section{Phase 2 results: the networks on Ibex}\label{sec:ibexres}

\subsection{One multiplication at a time}\label{sec:multbench}

First, the simplest measurement (\file{hw/sw/mult\_bench.c}, the paper's ``MULT kernel''): call each
routine 10{,}000 times on random unsigned numbers of 4, 8, 16 and 31 bits, and divide the cycle count by
10{,}000. This includes the call, loading the two numbers and the loop, like a real program.

\begin{figure}[H]
\centering
\includegraphics[width=\linewidth]{mult_bench}
\caption{Cycles per multiplication. Left: with the LOD instruction. Right: with LOD done in software.
The exact shift-and-add (black) grows with the size of the numbers. RoBA-Dn1 (blue, solid) stays at 28
cycles for any size, but only with the instruction: in software it costs 91--103 cycles.}
\end{figure}

\begin{table}[H]
\centering
\small
\begin{tabular}{lrrrr}
\toprule
 & \multicolumn{2}{c}{with LOD instruction} & \multicolumn{2}{c}{LOD in software}\\
\cmidrule(lr){2-3}\cmidrule(lr){4-5}
Multiplier & 8-bit & 31-bit & 8-bit & 31-bit\\
\midrule
Exact (shift-and-add) & 84.1 & 279.6 & 84.1 & 279.6\\
Mitchell & 33.5 & 34.0 & 96.5 & 109.0\\
DRUM6 & 98.5 & 101.0 & 148.8 & 164.0\\
RoBA-Dn1 & \textbf{28.0} & \textbf{28.0} & 91.0 & 103.0\\
RoBA-UpDn & 42.9 & 43.0 & 112.4 & 124.5\\
RoBA-Dn2 & 41.4 & 42.0 & 149.7 & 176.9\\
RoBA-UpDnDn & 71.5 & 72.0 & 192.7 & 219.8\\
RoBA-(UpDn)2 & 83.9 & 86.0 & 209.3 & 238.5\\
\bottomrule
\end{tabular}
\caption{Cycles per multiplication (\file{results/ibex/mult\_bench.csv}). Exact needs no LOD, so its two
columns are the same.}
\end{table}

\begin{keyidea}
With LOD, Dn1 is $84/28 = \mathbf{3.0\times}$ faster than exact for 8-bit numbers and
$280/28 = \mathbf{10\times}$ for 31-bit numbers. Without the instruction, it is \emph{slower} than exact.
The new instruction is what makes approximate multiplication worth it.
\end{keyidea}

\subsection{Whole networks, one routine call per multiplication}

Now the real thing: a whole picture through the network, written the way the paper does it. Every product
calls a small wrapper, \code{mul\_s(a, w)}, which checks for zero, takes the absolute values, calls the
multiplier routine and puts the sign back.

\begin{figure}[H]
\centering
\includegraphics[width=\linewidth]{ibex_call}
\caption{Million cycles per picture (mean of the first 100 test pictures), one routine call per
multiplication.}
\end{figure}

\begin{table}[H]
\centering
\small
\begin{tabular}{lrrrr}
\toprule
Multiplier & MNIST MLP & MNIST LeNet-5 & Fashion MLP & Fashion LeNet-5\\
\midrule
Exact & 16.89 & 29.62 & 16.04 & 29.32\\
Mitchell & 10.59 & 19.17 & 9.87 & 19.16\\
DRUM6 & 22.61 & 34.23 & 22.94 & 35.08\\
RoBA-Dn1 & \textbf{9.25} (1.83$\times$) & \textbf{17.73} (1.67$\times$) & \textbf{8.75} (1.83$\times$) & \textbf{17.59} (1.67$\times$)\\
RoBA-UpDn & 11.91 & 21.44 & 11.96 & 21.79\\
RoBA-Dn2 & 11.93 & 21.03 & 11.34 & 20.98\\
RoBA-UpDnDn & 17.36 & 28.60 & 17.95 & 29.48\\
RoBA-(UpDn)2 & 19.96 & 31.67 & 20.32 & 32.42\\
\bottomrule
\end{tabular}
\caption{Million cycles per picture, one call per product (\file{results/ibex/dnn\_call\_summary.csv}).
In brackets: how many times faster than exact.}
\end{table}

Two surprises:
\begin{enumerate}
  \item Dn1 is 3$\times$ faster per multiplication, but only \textbf{1.7--1.8$\times$} faster per picture.
  \item DRUM6, UpDnDn and (UpDn)2 are \textbf{slower than exact}. Network numbers are only 7 bits plus a sign,
        so the exact shift-and-add loop runs at most 7 rounds. That is cheap enough that methods doing
        two rounds of RoBA, or a small exact multiply, cannot beat it.
\end{enumerate}
The next subsection explains the first surprise.

\subsection{Where does the time go?}

\file{hw/profile\_ibex.py} runs one picture with Ibex's \emph{instruction tracer} switched on. It writes one
line for every instruction executed (about 1~GB of text per picture!), with the address of the instruction
and how many cycles it took. Using the program's debug information, every address is mapped back to the C
function and line it came from, and the cycles are added up per category. The traced total matches the CPU's
cycle counter within 0.1\%.

\begin{figure}[H]
\centering
\includegraphics[width=\linewidth]{profile}
\caption{Share of the time of one Fashion-MNIST picture, with one call per multiplication.}
\label{fig:profile}
\end{figure}

For the exact MLP, the multiplier routine itself takes \textbf{64\%} of the time. Another 17\% is the
sign-and-zero wrapper around it, and the rest is loading the numbers, the loop counter and adding up. For
Dn1, the routine shrinks to \textbf{35\%}, and the wrapper and loop now cost more than the multiplication.

\begin{example}[: Amdahl's law, the limit of speeding up one part]
Say a job takes 100 seconds and 64 seconds of it are multiplication. Even if multiplication became
\emph{free}, the other 36 seconds stay. The best possible speed-up is $100/36 = 2.8\times$.
In general: if a part takes a share $f$ of the time and you make it $k$ times faster, the whole job becomes
\[
  \text{speed-up} = \frac{1}{(1-f) + f/k}.
\]
For Dn1: $f = 0.64$, $k = 3$: $1/(0.36 + 0.213) = 1.74\times$. Close to the measured 1.83$\times$. The lesson:
after making the multiplier fast, the rest of the loop becomes the bottleneck.
\end{example}

\subsection{Writing the loop more cleverly: the fused MAC}\label{sec:fused}

To shrink the ``rest'', we wrote the multiply-accumulate loop differently (default build of
\file{hw/sw/dnn.c}, with the routines from \file{csrc/roba\_core.h}). We made the same three changes for
\emph{every} multiplier, exact included, so the comparison stays fair.

\paragraph{1. Inline the routine.} Instead of calling a function for every product (jump there, save
registers, jump back), the multiplier's few instructions are copied directly into the loop. The zero check
is done once per product (the old wrapper and the routine each checked).

\paragraph{2. Store weights as sign and size.} Normally $-5$ is stored as \code{11111011} (``two's
complement''), and getting its size needs a test and a negation. We store each weight as one byte with
bit~7 = the sign and bits 0--6 = the size:
\[
  +5 \to \texttt{0\,0000101} = \texttt{0x05},\qquad -5 \to \texttt{1\,0000101} = \texttt{0x85}.
\]
Then the size is one instruction (\code{byte \& 127}) and the sign another (\code{byte >> 7}). This
conversion is done once, when the program starts, outside the timed part (a real device would store the
weights like this).

\paragraph{The sign without a jump.} After multiplying the sizes, the result $r$ must be negated if the two
signs differ. A jump (``if negative, negate'') costs extra cycles on Ibex when it is taken. Instead we use a
trick: let $s = 0$ if the signs agree and $s = -1$ (all bits 1) if they differ. Then
\[
  (r \oplus s) - s = \begin{cases} r & s = 0\\ -r & s = -1\end{cases}
\]
because XOR with all ones flips every bit ($\sim r = -r-1$), and subtracting $-1$ adds 1.
\begin{example}[: the branch-free sign]
$r = 128$, $s = -1$: $128 \oplus (-1) = -129$, and $-129 - (-1) = -128$. Correct, with no jump.
With $s=0$: $128\oplus 0 - 0 = 128$.
\end{example}

\paragraph{3. After ReLU, inputs are never negative.} Every layer except the first gets its input from a
ReLU, so the input is $\ge 0$: no need to take its absolute value or look at its sign.

\begin{lstlisting}[language=C]
/* acc + a*w for an activation known to be >= 0 (after ReLU), from hw/sw/dnn.c */
static inline int32_t mac_r(int32_t acc, int32_t a, uint8_t w)
{
    uint32_t ub = w & 0x7Fu;                    /* |w|                         */
    if (a == 0 || ub == 0)                      /* one zero check              */
        return acc;
    int32_t r = (int32_t)MULT_NZ((uint32_t)a, ub);   /* e.g. roba_dn1_nz, inlined */
    int32_t s = -(int32_t)((uint32_t)w >> 7);   /* -1 if w < 0, else 0         */
    return acc + ((r ^ s) - s);                 /* add r or -r, no branch      */
}
\end{lstlisting}

\begin{figure}[H]
\centering
\includegraphics[width=\linewidth]{fused}
\caption{Fashion-MNIST: one call per multiply (white circles) against fused into the loop (filled). Every
multiplier gets faster; RoBA-Dn1 gains the most on LeNet-5.}
\end{figure}

\begin{table}[H]
\centering
\small
\begin{tabular}{lrrrr}
\toprule
Multiplier (fused) & MNIST MLP & MNIST LeNet-5 & Fashion MLP & Fashion LeNet-5\\
\midrule
Exact & 13.91 & 24.04 & 13.49 & 22.95\\
Mitchell & 8.61 (1.62$\times$) & 11.82 (2.03$\times$) & 8.14 (1.66$\times$) & 12.16 (1.89$\times$)\\
DRUM6 & 15.87 & 26.01 & 15.45 & 27.18\\
RoBA-Dn1 & \textbf{7.12 (1.95$\times$)} & \textbf{10.19 (2.36$\times$)} & \textbf{7.05 (1.91$\times$)} & \textbf{10.21 (2.25$\times$)}\\
RoBA-UpDn & 10.13 & 15.42 & 9.73 & 15.70\\
RoBA-Dn2 & 10.00 & 13.62 & 9.85 & 13.78\\
RoBA-UpDnDn & 14.94 & 22.33 & 14.42 & 23.03\\
RoBA-(UpDn)2 & 17.87 & 25.87 & 17.25 & 26.78\\
\bottomrule
\end{tabular}
\caption{Million cycles per picture with the fused loop (\file{results/ibex/dnn\_summary.csv}); in brackets
the speed-up over fused exact. Compared with the call version: exact 16.0 $\to$ 13.5 (Fashion MLP), Dn1
8.75 $\to$ 7.05, and Dn1 on Fashion LeNet-5 17.6 $\to$ 10.2 ($1.72\times$).}
\end{table}

Here is what the compiler made of the fused Dn1 loop for the MLP's first layer: 26 instructions per product,
and only 10 when the input is 0 (it jumps straight to label 2).
\begin{lstlisting}[language=asm]
      lb    a5,0(a0)      # load x
      addi  a0,a0,1
      j     1f
1:    srai  a7,a5,31      # sign of x
      xor   t3,a7,a5
      lbu   a3,0(a1)      # load w (sign-magnitude byte)
      sub   t3,t3,a7      # |x|
      beqz  a5,2f         # x = 0: skip
      andi  a5,a3,127     # |w|
      beqz  a5,2f         # w = 0: skip
      lod   a6,t3         # pa = LOD(|x|)      <- the new instruction
      lod   t1,a5         # pb = LOD(|w|)      <- the new instruction
      sll   a5,a5,a6      # |w| << pa
      srli  a4,a3,7       # sign of w
      sll   t3,t3,t1      # |x| << pb
      sll   a6,t5,a6      # 1 << pa
      neg   a4,a4
      add   a5,a5,t3
      sll   a6,a6,t1      # 2^(pa+pb)
      xor   a4,a4,a7      # combined sign s
      sub   a5,a5,a6      # Dn1 product r
      xor   a5,a5,a4      # (r ^ s)
      sub   a5,a5,a4      #   - s
      add   a2,a2,a5      # acc += product
2:    addi  a1,a1,1
      bne   a1,t0,loop
\end{lstlisting}

\subsection{No multiplier at all: a table of squares}\label{sec:sqtab}

There is an old trick (it was used by people multiplying with tables of squares before calculators). Since
$(x+w)^2 = x^2 + 2xw + w^2$,
\[
  x\cdot w = \frac{(x+w)^2 - x^2 - w^2}{2}.
\]
A square can be \emph{looked up} in a table instead of computed. So a product becomes three lookups, two
subtractions and a halving. This is called the \textbf{polarization identity}.

\begin{example}[: products from squares]
$3\times5$: $\big((3+5)^2 - 3^2 - 5^2\big)/2 = (64 - 9 - 25)/2 = 30/2 = 15$.\\
$3\times(-5)$: $\big((-2)^2 - 9 - 25\big)/2 = (4-34)/2 = -15$. Negative numbers just work.
\end{example}

It gets much better for a whole neuron. Add up the identity over all inputs:
\[
  \sum_k x_k w_k = \frac{1}{2}\Big(\underbrace{\sum_k S(x_k+w_k)}_{\text{1 lookup per MAC}}
  - \underbrace{\sum_k S(x_k)}_{\text{same for every neuron}}
  - \underbrace{\sum_k S(w_k)}_{\text{fixed: computed in advance}}\Big),\qquad S(v) = v^2.
\]
\begin{itemize}
  \item $\sum S(w_k)$ depends only on the weights, so it is computed once, before running (outside the
        timed part, like the weights themselves).
  \item $\sum S(x_k)$ depends only on the layer's input, which is the \emph{same} for all 256 neurons of the
        MLP's first layer. It is computed once and shared.
  \item So each MAC costs about \textbf{one table lookup and one addition}.
\end{itemize}

\begin{example}[: a two-input neuron]
$x = (3, 2)$, $w = (5, -1)$. Exact: $15 - 2 = 13$.
$\sum S(x+w) = S(8) + S(1) = 64 + 1 = 65$; $\sum S(x) = 9 + 4 = 13$; $\sum S(w) = 25 + 1 = 26$.
$(65 - 13 - 26)/2 = 26/2 = 13$. Exact, with no multiplication. (The number before halving is always
even, because it equals $2\sum x_kw_k$, so the halving is exact.)
\end{example}

\paragraph{The table.} Inputs and weights are int8 ($-127\dots127$), so $x+w$ is in $-254\dots254$: 509
entries. The largest square, $254^2 = 64{,}516$, fits in 16 bits, so the table is $509\times 2 =
1{,}018$ bytes, about \textbf{1~KB}. It is built at start-up without any multiplication, using
$(n+1)^2 = n^2 + 2n + 1$: start at 0 and add $1, 3, 5, 7, \dots$ $\to 0, 1, 4, 9, 16, \dots$

\paragraph{Convolutions.} In a convolution each output position sees a different $5\times5$ window of
the input, so $\sum S(x)$ is different per position, but it is still the same for all filters at that
position. The program computes all window sums once per layer with ``box sums'' (add 5 neighbours across,
then 5 down) and shares them across the 6 (or 16) filters.

The inner loop is now 10 instructions instead of 26, and none of them is LOD:
\begin{lstlisting}[language=asm]
      lb    a5,0(a4)      # load x
      lb    a0,0(a3)      # load w
      addi  a4,a4,1
      addi  a3,a3,1
      add   a5,a5,a0      # x + w
      slli  a5,a5,1       # times 2 bytes per entry
      add   a5,a5,a1      # address in the table
      lhu   a5,0(a5)      # S[x + w]
      add   a2,a2,a5      # d += S[x + w]
      bne   s1,a4,loop
\end{lstlisting}

\begin{figure}[H]
\centering
\includegraphics[width=\linewidth]{sqtab}
\caption{Million cycles per picture and speed-up over fused exact. The table of squares is the fastest
everywhere, and its products are exact.}
\end{figure}

\begin{figure}[H]
\centering
\includegraphics[width=\linewidth]{layers}
\caption{Cycles per layer (Fashion-MNIST). In the MLP almost everything is the first layer
($784\to256$, 85\% of the multiplications). In LeNet-5, conv2 is the largest part.}
\end{figure}

\begin{table}[H]
\centering
\small
\begin{tabular}{lrrrr}
\toprule
 & MNIST MLP & MNIST LeNet-5 & Fashion MLP & Fashion LeNet-5\\
\midrule
Exact, fused & 13.91 & 24.04 & 13.49 & 22.95\\
RoBA-Dn1, fused & 7.12 (1.95$\times$) & 10.19 (2.36$\times$) & 7.05 (1.91$\times$) & 10.21 (2.25$\times$)\\
Square table & \textbf{3.87 (3.59$\times$)} & \textbf{7.58 (3.17$\times$)} & \textbf{3.87 (3.49$\times$)} & \textbf{7.57 (3.03$\times$)}\\
\bottomrule
\end{tabular}
\caption{Million cycles per picture. The table is 1.3--1.8$\times$ faster than fused Dn1, with exact
products and no LOD.}
\end{table}

\begin{kid}
Notice that the table method takes \emph{exactly} the same time (3.87 million cycles) for an MNIST
digit and a Fashion-MNIST shoe. It does the same lookups no matter what the numbers are. The RoBA and
exact loops skip zeros and their speed depends on the numbers, so their times differ a little between
datasets.
\end{kid}

\begin{careful}
The table trick depends on two things that are true here but not everywhere. (1) Ibex's memory answers in
one cycle; with slow flash memory or a cache, every lookup could cost more. (2) The numbers are 8-bit. For
16-bit numbers the table would need about 131{,}000 entries of 4 bytes each, roughly \textbf{512~KB}, too
big for a small chip, while the LOD-based RoBA routines cost the same for any size.
\end{careful}

\paragraph{A tempting idea that does not work: approximate squares.} One could replace the table with
RoBA-style \emph{approximate} squares, $\tilde S(x) = (|x| \ll (k+1)) - (1\ll 2k)$ with $k = \text{LOD}(|x|)$.
This is a straight-line estimate of $x^2$ that is always a little too small. The problem is that the
three squares are big numbers, and the product is the small difference between them, so small errors in
the squares become large errors in the product:
\begin{example}[: why approximate squares fail]
$3\times4 = 12$. $\tilde S(7) = (7\ll3) - (1\ll4) = 56-16 = 40$ (true 49); $\tilde S(3) = 12-4 = 8$
(true 9); $\tilde S(4) = 16$ (exact, a power of two).
$(40 - 8 - 16)/2 = 8$ instead of 12: \textbf{33\% too small}, although no square was more than 18\% off.
\end{example}
We tested it on the GPU: the networks lost 6--15 accuracy points in 3 of the 4 cases, so it was rejected.
The exact table is both exact and fastest.

\subsection{Skipping zeros, and feeding raw pixels}\label{sec:zeroskip}

Anything times zero is zero, so a zero input could be skipped entirely. How many zeros are there?
\begin{itemize}
  \item In \textbf{hidden layers}, many: ReLU turns every negative value into 0.
  \item In the \textbf{input picture}, none! Remember that normalising turns black into $-0.424$ (MNIST),
        not 0. Most of an MNIST picture is black background, all of it non-zero.
\end{itemize}

\paragraph{Raw pixels: folding the normalisation into the first layer.} The first layer computes
$\sum_k w_k\,\frac{p_k - \text{mean}}{\text{std}} + b$. This can be rewritten as
\[
  \sum_k \underbrace{\frac{w_k}{\text{std}}}_{\text{new weight}}\, p_k
  \;+\; \underbrace{\Big(b - \text{mean}\sum_k \frac{w_k}{\text{std}}\Big)}_{\text{new bias}}.
\]
Same result, but now the network reads the raw pixel $p_k$, and black is exactly 0.
\begin{example}[: folding one weight]
MNIST, $w = 0.6$, pixel $p = 0.5$. Original: $0.6\times(0.5-0.1307)/0.3081 = 0.6\times1.1986 = 0.719$.
Folded: new weight $0.6/0.3081 = 1.947$, bias change $-0.1307\times1.947 = -0.254$;
$1.947\times0.5 - 0.254 = 0.719$. The same.
\end{example}
The raw int8 input is $\text{round}(127\,p)$, computed from the original 8-bit pixel $k$ as
$(254k + 255)\,\text{div}\,510 = \text{round}(127k/255)$. After this, \textbf{81\% of MNIST pixels and 51\%
of Fashion-MNIST pixels are exactly 0}.

\begin{careful}
For LeNet-5 there is a detail at the border. conv1 pads the picture with zeros so that the output is still
$28\times28$. With normalised inputs, a padded 0 means ``an average pixel''; with raw inputs it means
``black''. So the new bias depends on how many of the filter's 25 taps fall on the real picture, and conv1
gets one bias per output position (a $6\times28\times28$ map) instead of one per filter
(\code{C0\_PER\_POSITION} in the code).
\end{careful}

\paragraph{Skipping (``+zs'').} The multiplier loops already skip zero products one at a time (the zero
check). With \code{ZERO\_SKIP}, a fully connected layer first collects the non-zero inputs and their positions
once, and then every neuron loops only over those. For the table method skipping needs more work per
remaining term: a zero input adds $S(0+w) - S(0) - S(w) = 0$, so it can be dropped, but then the precomputed
$\sum S(w)$ no longer matches, and the loop must use $S(x+w) - S(w)$: \emph{two} lookups per term.

\begin{figure}[H]
\centering
\includegraphics[width=\linewidth]{zeroskip}
\caption{The ablation: each method with normalised or raw pixels, with or without skipping. The labels give
the speed-up over exact with normalised input (the first bar).}
\end{figure}

\begin{table}[H]
\centering
\small
\begin{tabular}{llrrrr}
\toprule
Method & Setting & MNIST MLP & MNIST LeNet-5 & Fashion MLP & Fashion LeNet-5\\
\midrule
\multirow{4}{*}{Exact} & normalised & 13.91 & 24.04 & 13.49 & 22.95\\
 & + skip & 13.49 & 23.85 & 13.96 & 22.76\\
 & raw & 5.45 & 17.81 & 7.75 & 19.73\\
 & raw + skip & 3.18 & 17.64 & 6.82 & 19.60\\
\midrule
\multirow{4}{*}{RoBA-Dn1} & normalised & 7.12 & 10.19 & 7.05 & 10.21\\
 & + skip & 6.92 & 9.94 & 7.00 & 10.16\\
 & raw & 3.57 & 8.86 & 4.30 & 9.68\\
 & raw + skip & 1.58 & 8.69 & 3.23 & 9.62\\
\midrule
\multirow{4}{*}{Square table} & normalised & 3.87 & 7.58 & 3.87 & 7.57\\
 & + skip & 5.10 & 8.20 & 4.99 & 8.31\\
 & raw & 3.87 & 7.57 & 3.64 & 7.62\\
 & raw + skip & \textbf{1.33} & \textbf{7.38} & \textbf{2.70} & 7.93\\
\bottomrule
\end{tabular}
\caption{Million cycles per picture (\file{results/ibex/zeroskip\_summary.csv}). Best per column in bold
(Fashion LeNet-5: the table with raw pixels and no skip, 7.62, is about as fast as the normalised table).}
\end{table}

What this shows:
\begin{itemize}
  \item \textbf{Raw pixels are the big win.} Exact on the MNIST MLP drops from 13.9 to 5.5 million cycles just
        by giving the existing zero check something to skip.
  \item \textbf{The best MLP result:} table + raw + skip on MNIST, \textbf{1.33 million cycles, $10.5\times$ faster}
        than exact with normalised input. On Fashion (fewer black pixels) it is $5.0\times$.
  \item \textbf{Skipping without raw pixels does not help}, and it makes the table \emph{slower}
        ($3.87\to4.99$ on Fashion), because each remaining term now needs two lookups.
  \item \textbf{LeNet-5 hardly changes}: most of its work is conv2, whose input comes after ReLU and pooling
        and is not affected by the raw pixels, and the convolution can only skip tap by tap.
\end{itemize}

\begin{table}[H]
\centering
\small
\begin{tabular}{lcccc}
\toprule
 & \multicolumn{2}{c}{MNIST MLP} & \multicolumn{2}{c}{Fashion MLP}\\
\cmidrule(lr){2-3}\cmidrule(lr){4-5}
 & normalised & raw & normalised & raw\\
\midrule
Exact / table & 98.41 & 98.46 & 89.75 & 89.69\\
RoBA-Dn1 & 98.36 & 98.42 & 89.53 & 89.53\\
\bottomrule
\end{tabular}
\caption{Accuracy on all 10{,}000 test pictures (from the bit-exact Python model) with normalised and raw
input: raw pixels cost no accuracy. Skipping never changes a result. (LeNet-5: MNIST 99.18/99.19, Fashion
90.83/90.80 for exact.)}
\end{table}

\subsection{Energy}\label{sec:energy}

Energy is power multiplied by time: $E = P\times T$. Time is proportional to cycles. For power we use the
paper's synthesis results, relative to a core with a single-cycle hardware multiplier ($P = 1$):
\begin{itemize}
  \item Multiplierless core \emph{with} LOD: 17.1\% less power, so $P_{\text{LOD}} = 1 - 0.171 = 0.829$.
  \item LOD itself adds 4.5\%, so the multiplierless core \emph{without} LOD has
        $P_{\text{none}} = 0.829/1.045 = 0.793$.
\end{itemize}
Exact and the table of squares do not use LOD, so they run on the cheaper 0.793 core. Mitchell, DRUM6 and all
RoBA methods need the LOD core (0.829).

\begin{example}[: energy of one Fashion-MNIST MLP picture (fused)]
\begin{tabular}{@{}lrl@{}}
Exact: & $13.49 \times 0.793 = 10.70$ & (reference, 1.00)\\
RoBA-Dn1: & $7.05 \times 0.829 = 5.84$ & $\Rightarrow 5.84/10.70 = \mathbf{0.55}$\\
Square table: & $3.87 \times 0.793 = 3.07$ & $\Rightarrow \mathbf{0.29}$\\
RoBA-(UpDn)2: & $17.25\times0.829 = 14.30$ & $\Rightarrow 1.34$ (more energy than exact!)
\end{tabular}
\end{example}

\begin{figure}[H]
\centering
\includegraphics[width=\linewidth]{energy}
\caption{Energy per picture relative to exact (fused, Fashion-MNIST). Below the dashed line saves energy.}
\end{figure}

\begin{keyidea}
Because LOD raises the power by 4.5\%, an approximate method has to be at least 4.5\% faster than exact
just to break even. Dn1 is about 2$\times$ faster, so it needs only \textbf{0.44--0.55} of the energy.
The two-step methods and DRUM6 are slower than exact and \emph{waste} energy. The table of squares needs
\textbf{0.28--0.33}.
\end{keyidea}

\begin{careful}
These energies are \emph{estimates}: measured cycles times the paper's relative power numbers. The real
power also depends on what the program does (memory accesses use more than additions, for example). A
real power measurement needs the FPGA or a chip.
\end{careful}

\subsection{How do we know the simulated chip computed the right thing?}

\begin{itemize}
  \item \textbf{All 10{,}400 different inferences} (100 pictures $\times$ 104 distinct configurations) gave
        \textbf{exactly the same scores, bit for bit,} as the Python model of the same integer arithmetic
        (\code{logits\_match\_ref = 1.0} in every row of the CSV files).
  \item The 32 ``call'' configurations were run in two separate runs, from rebuilt programs, and gave
        identical cycle counts and answers.
  \item Ibex accuracy on these 100 pictures: MNIST 100\% for every configuration; Fashion 90--91\% (MLP)
        and 85--88\% (LeNet-5). With 100 pictures one picture is one point, so these numbers move more than
        the full-test numbers. The full-test accuracy comes from the bit-exact Python model.
  \item The integer Ibex model agrees with the GPU's int8 model on 97--100\% of the 100 pictures. The small
        difference is not an error: the GPU version rescales with decimal arithmetic, the integer version with
        $M$ and a shift, so a value right at a rounding edge (like 6.5) can round the other way and, rarely,
        flip an answer. Over the full test set the two agree within 0.05 points for exact and within about
        0.3 points for RoBA-Dn1 (Fashion LeNet-5: 90.49\% integer vs 90.77\% GPU).
\end{itemize}
\newpage
